\chapter{Discussion}\label{chap:discussion}

This chapter interprets the findings of Chapter~\ref{chap:experiments}, puts them next to the prior work from Chapter~\ref{chap:background}, and sets out what the study cannot claim. It is organised around the two research questions, the lag structure of the news-liquidity response (RQ1) and its directional asymmetry (RQ2), followed by the methodological contributions of the two-phase design and the limitations that qualify everything above them.

\section{Lag structure of the news-liquidity response (RQ1)}\label{sec:rq1-discussion}

RQ1 asks at what lag news events hit WTI liquidity hardest. Two separate analyses land on the same answer.

\begin{itemize}
  \item \textbf{Phase~1's lag OLS} (Section~\ref{sec:lag-ols}) places the peak of the bearish news effect on \texttt{log\_volume} at $+6$ hours, with the effect built up over the preceding hours and decayed to insignificance by $+8$ hours.
  \item \textbf{Phase~2's TFT v2} corroborates this (Section~\ref{sec:tft-v2-lag}) through its encoder attention, which, disaggregated by sentiment direction, peaks at $-6$ hours for bearish-sentiment hours, aligning exactly with the lag OLS peak. Its per-horizon error curve is a weaker instrument for this purpose: it confirms that the Phase~2 features carry predictive signal across the whole 1-to-12 hour window, but it rises mainly because the persistence baseline decays with horizon, so it does not independently locate the response at any lag.
\end{itemize}

One of these is a regression coefficient estimated on a single sentiment value at a fixed lag. The other is a deep-learning forecaster integrating a 48-hour window of many features at once. They have almost nothing in common, and they pick out the same lag, \emph{$+6$ hours}. That agreement, on a point rather than a broad window, is the strongest evidence this thesis has for RQ1.

What it means is that news does not reprice liquidity on contact. It propagates over several hours. If the response were purely contemporaneous the peak would sit at lag zero, and if the market absorbed news fully and immediately there would be no predictable structure left six hours out to find.

What we see instead is a response that builds and then persists across the session and into the next hour. That fits gradual information diffusion and traders reacting at different times, rather than one clean repricing event. The TFT attention fills the picture in: a strong recency component at $-1$ hour, since the most recent news is always the most useful thing for an immediate forecast, sitting alongside a longer integration window that applies to bearish news specifically.

So the market seems to \textbf{take longer} to work through news implying a \emph{downward} price move than news implying an \emph{upward} one. That is a directional detail inside the temporal structure. Phase~1's average-lag estimate cannot see it, and the TFT's direction-conditioned attention can.

All of this is consistent with the premise from the literature (Section~\ref{sec:related-work}) that news sentiment moves trading activity and not only prices \citep{tetlock2007}. What is new here is resolution. Work at \textbf{daily} frequency cannot see an intraday six-hour propagation window at all, and the \textbf{hourly} design of this thesis exists to catch exactly that. The literature already establishes that news effects concentrate at high frequency, mostly within hours of release \citep{ederington1993}. Practically, the result says that an hourly monitor of WTI-relevant news carries liquidity information that is most actionable not in the same hour, but roughly six hours later, within the same trading day.

\section{Directional asymmetry of the response (RQ2)}\label{sec:rq2-discussion}

RQ2 asks whether bearish news produces a different liquidity response than comparable bullish news. Phase~1's lag OLS (Section~\ref{sec:lag-ols}) says yes: the coefficient on bearish sentiment beats the coefficient on bullish sentiment at the dominant lags, and that ordering is the thesis' primary evidence for RQ2. Two qualifications belong with it from the outset.

The first is what kind of statement it is. It is about marginal sensitivity, meaning how much volume moves for an increment of bearish rather than bullish sentiment, with the other held fixed. The second is how strong it is. Each coefficient is individually significant at the peak ($p < 0.001$) and bearish exceeds bullish at every dominant lag, by about 18\% at $+6$h, but the difference between the two coefficients is not itself submitted to an equality test. The evidence is therefore a consistent ordering across lags, reinforced by an independent method in Phase~2's attention, rather than a single tested contrast. A formal test of $\beta_1 = \beta_2$ would sharpen this, and Section~\ref{sec:future-work} returns to it.

Phase~2's TFT v2 does not reproduce the asymmetry in its point predictions. Across sixteen bearish-versus-bullish comparisons of predicted volume (Section~\ref{sec:asymmetry}), none is significant. That looks like a failed replication and it is not one, because the two analyses are estimating different quantities.

The lag OLS compares marginal coefficients, which are slopes with other variables held fixed. The TFT test compares the average predicted volume of bearish hours against bullish hours, which is an unconditional difference of group means with nothing held fixed. Those can both be true at once. A larger marginal sensitivity to bearish news can sit alongside near-equal average volume across bearish and bullish hours, for example if bullish hours happen to fall in periods where baseline volume is already high. The null is also not a power artifact. The two groups each hold hundreds of samples, the result survives sweeping the sentiment threshold from 0.0 to 0.5, and the target itself is near-symmetric: in the actual test data, future volume for bearish and bullish hours is statistically indistinguishable at the 1 to 6 hour horizons. The model even reproduces the one real difference, a bullish-above-bearish gap at $+12$ hours. It is tracking a near-symmetric ground truth rather than washing out a signal that was there.

There is a second reason the phases point different ways, and it runs deeper: they do not define sentiment the same way (Section~\ref{sec:phase-comparison}). FinBERT scores the surface financial tone of an article. The Haiku \texttt{sentiment\_score} is prompted for the net directional impact on the WTI price. The two agree only moderately (Pearson 0.54), and they disagree in sign on exactly one kind of news, which is geopolitical events.

\begin{quote}
A headline such as ``Russian Sanctions Drive Oil Prices Higher'' is negative in tone (sanctions, conflict) but bullish in price (a supply threat lifts crude), so FinBERT labels it bearish while Haiku labels it bullish.
\end{quote}

That geopolitical news is the high-volume, high-risk-premium subset that dominates the liquidity response in the first place. Under Phase~1's tone-based definition it lands in the bearish bucket, which reinforces the bearish-over-bullish asymmetry. Under Phase~2's price-based definition it lands in the bullish bucket, which is why the TFT's price-conditioned test finds no bearish-over-bullish gap and even a mild reversal. So the tension between the phases is not a contradiction. It is what happens when bearishness is measured two different ways.

\textbf{For the TFT it is not a matter of negative versus positive news, but of the \emph{risk premium} the news implies.}

Looking at what the TFT actually leans on makes this clearer. Its most important features (Section~\ref{sec:feature-importance}) are VIX, the supply and demand channels, and the entity block, whose 71 flags carry 52\% of total importance between them while the composite \texttt{sentiment\_score} sits 21st. None of those is a bearish signal in the price sense. VIX is a volatility expectation, the channels are fundamental supply and demand judgments, and an entity flag records only that a story named a given actor, carrying no direction at all. What the model weights, in other words, is how much a story raises supply uncertainty and how salient its subject is, not which way it points. Which specific entities rank highest is not stable enough to interpret individually (Section~\ref{sec:feature-importance}), so the argument rests on the block and on the feature types, not on a roster of actors. So the asymmetry the model learned is not bearish over bullish. It is risk-and-salience over baseline: WTI liquidity responds to how much a piece of news raises uncertainty about supply, not to whether it implies a higher or lower price.

The answer to RQ2 therefore depends on how sentiment is defined and which quantity is being estimated. There is a robust bearish-over-bullish asymmetry in the marginal sensitivity of volume to negative-tone news (Phase~1), and that stays the thesis' primary answer. It does not show up as a difference in the level of predicted volume between price-bearish and price-bullish hours (Phase~2), because the volume-driving geopolitical news is price-bullish, and because the model organises its predictions around risk and salience rather than price direction. Both readings are right inside their own framing. Stating the asymmetry properly means naming the sentiment construct and the estimand along with it.

\section{Methodological contributions}\label{sec:methodological-contributions}

Beyond the two research questions, the thesis makes five methodological contributions. Each one came out of a concrete problem we ran into, not from deciding in advance to contribute something.

\textbf{Aligning a continuous news stream onto a discontinuous trading grid.} Sub-daily work creates a problem that daily work does not have. News publishes around the clock; the market does not trade around the clock. So before any hourly analysis can begin, every article has to be placed on a grid punctuated by overnight, weekend and holiday closures. Two rules do that here (Section~\ref{sec:temporal-alignment}). \emph{Ceiling assignment} sends an article published at 14:23 to the 15:00 bar rather than the 14:00 one, because the 14:00 bar had already been forming for twenty-three minutes when the article appeared, and letting news explain a candle that predates it would break the causal reading outright. \emph{Forward-assignment} then carries off-hours articles to the next tradeable hour instead of discarding them, which matters because weekend OPEC communiqu\'es and overnight geopolitical events are exactly the highest-impact items in the corpus. Both rules are conservative in the same direction, understating the immediate reaction rather than inflating it, so the $+6$h peak is if anything a lower bound on how quickly the market responds. Setting aside the boundary case below, 22.9\% of articles are forward-assigned, with a median gap of half an hour and a maximum of 73.8 hours across a holiday weekend. That is the scale of the problem being solved rather than a marginal correction, and it is a prerequisite for intraday work that a daily design never has to confront.

\emph{One boundary defect belongs in this account.} Forward-assignment has no natural answer for articles published \emph{before} the market grid begins, and the implementation carries all of them to the first available hour: 3{,}657 articles, 16\% of the corpus, collapse onto 13 May 2024 12:00 with a mean assignment gap of 1{,}617 hours. That hour is index 1 of the grid and sits inside the training partition, so it never enters validation or test and reaches only the encoder window of the earliest training samples; its effect on the reported results is negligible. The correct handling is to drop pre-grid articles rather than forward-assign them, and any hourly corpus statistic evaluated at that timestamp should be read as an artefact rather than as data.

\textbf{Channel decomposition as a way to calibrate model annotation.} Splitting sentiment into orthogonal supply, demand, and risk-premium channels (Section~\ref{sec:channel-decomposition}), following the structural decomposition of oil-price shocks \citep{kilian2009}, was a fix for a specific failure: the composite \texttt{sentiment\_score} disagreed across model families (Haiku versus GPT-5.5) on high-magnitude geopolitical events (Section~\ref{sec:calibration}). Recalibrating on the channels lifted cross-model agreement on the composite from 0.39 to 0.88, with channel correlations of 0.94, 0.96, and 0.82. So the channels do two jobs: they are higher-resolution features for the TFT, and with no human ground truth available, cross-model agreement on them is a defensible validation signal for the extraction \citep{gilardi2023, zheng2023}.

\textbf{Tone versus price impact.} A tone classifier and a price-reasoning LLM are not a coarse and a fine version of the same measurement. They measure different things (Section~\ref{sec:phase-comparison}, Section~\ref{sec:rq2-discussion}). FinBERT reports an article's surface financial tone. The LLM does the asset-specific causal step, from a supply threat to a higher crude price, which a tone classifier has no way to do. The geopolitical sign-flip is the demonstration. The lesson travels beyond this thesis: for an asset-specific task, sentiment has to be defined relative to the asset and its causal mechanism rather than as generic positive-or-negative tone, and a study that mixes the two can reach opposite conclusions depending on which one it happened to pick.

\textbf{LLM-as-filter for noisy news corpora.} Phase~1 decided what to include with a brittle regex heuristic over the article body (Section~\ref{sec:regex-filter}). Phase~2 replaced it with the LLM's own judgment of material relevance (\texttt{usable}). The regex is lexical and the LLM filter is semantic, and comparing them (Section~\ref{sec:filter-comparison}) shows where each one breaks. The training filter, \texttt{usable\_strict}, then adds a non-zero channel score on top, but it is worth keeping the two ideas apart: \texttt{usable} is the relevance judgment and the actual contribution here, while \texttt{usable\_strict} is a signal requirement that drops on-topic articles which simply carry no channel information (Section~\ref{sec:filter-comparison}). For web-scraped news, where relevance is about meaning rather than keywords, a semantic filter is a practical alternative to keyword heuristics and transfers to other corpora.

\textbf{Phase-based progression as a reproducible template.} The two-phase design is a contribution in itself. A simple interpretable baseline (FinBERT features with OLS) answers the research questions with classical statistics whose assumptions you can see. A richer extraction with a more expressive model (LLM channels with a TFT) then corroborates the baseline and adds vocabulary to the analysis. Neither phase gets discarded. The second does not supersede the first, it triangulates it, and both have their own limitations. This is a reusable template for empirical news-impact studies that want interpretability and expressiveness without giving up either one.

\section{Limitations}\label{sec:limitations}

Several limitations qualify the conclusions and mark out how far they reach.

\textbf{Inter-model calibration is weaker than human ground truth.} We validate the LLM features by cross-model agreement (Haiku versus GPT-5.5, Section~\ref{sec:calibration}) rather than against expert human annotation, which the project could not afford (Section~\ref{sec:calibration-methods}). Cross-family agreement tells us something real, since two models built by different developers converging on a reading is not nothing, but it is not a gold standard. Both models can share a blind spot, and agreeing is not the same as being right. The calibration sample is also small (\emph{$n = 30$}) by design, which is enough to surface systematic disagreements and not enough to estimate population agreement precisely.

\textbf{Regime extrapolation on \texttt{price\_range}.} The most concrete modeling limitation is TFT v2 failing on \texttt{price\_range} in the war regime (Section~\ref{sec:predictive-performance}). Trained only on the pre-war period, the model falls back toward the historical mean when asked about a high-volatility regime it never saw, and it loses to a naive persistence baseline. This is a clean illustration of a model constrained by what its training window contained. It is also a warning not to assume the aggregate liquidity results (\texttt{log\_volume}, \texttt{amihud}), which do transfer across the regime boundary, will generalize to every target.

\textbf{Construct dependence of the sentiment findings.} As Section~\ref{sec:rq2-discussion} spells out, the directional conclusions depend on how sentiment is defined (tone versus price impact) and which quantity is estimated (marginal coefficient versus group mean). Stated carefully this is a strength, but it is a limitation for anyone who wants a single unconditional answer. There is no framing-free version of the asymmetry, and any restatement has to carry its definitions with it.

\textbf{Corpus composition.} The news corpus comes from GDELT (Section~\ref{sec:news-data}), a broad machine-monitored aggregator that skews toward retail-facing web and print coverage rather than the institutional feeds, trading terminals and wire services, that professional oil traders actually read. So what the models extract is the sentiment of publicly available news, which may lead, lag, be biased, or just differ from the information set that moves institutional order flow.

\textbf{The Phase~1 lag regressions cannot be re-fit from the archived data.} Hourly price history is served by \texttt{yfinance} on a rolling window of roughly two years, so the market grid is not a fixed asset. Phase~1's grid began on 11 March 2024, and by the time the data were re-pulled for Phase~2 the earliest hours had aged out; the retained hourly file now starts on 13 May 2024 (Section~\ref{sec:phase1-data}). Around 1{,}700 Phase~1 articles fall before that boundary, so the lag regressions of Section~\ref{sec:lag-ols} cannot be reproduced from the repository as it stands, and the archived coefficient table is the surviving record of them. The Phase~2 results are unaffected, since they run on the retained grid. Anyone repeating this design should snapshot the price grid alongside the derived features rather than assume it can be re-fetched.

\textbf{Non-stationarity and single-commodity scope.} The dataset covers roughly 24 months and contains a structural break, the Iran war onset (Section~\ref{sec:phase2-data}), which makes the series non-stationary across the train/test boundary. The war-regime results come from one episode and should be read as suggestive rather than established. The study also covers one instrument, WTI crude oil futures. The channel decomposition and the risk-and-salience reading of the liquidity response are plausibly specific to a commodity with WTI's geopolitical supply structure, and neither the lag nor the asymmetry finding can be assumed to carry to other commodities or asset classes without someone replicating it.

\section{Future work}\label{sec:future-work}

The limitations point to several extensions, and most of them are small enough to be tractable.

\textbf{A formal test of the directional asymmetry.} RQ2's primary evidence is an ordering of two coefficients that are each significant, not a tested difference between them (Section~\ref{sec:rq2-discussion}). The direct remedy is cheap: a Wald test of $\beta_1 = \beta_2$ on the $+6$h regression, which the fitted model already supports since it needs only the estimated covariance of the two coefficients, reported alongside a confidence interval for $\beta_1 - \beta_2$. Bootstrapping the difference across lags would additionally show whether the ordering is stable or driven by the peak alone. This would convert the asymmetry from a consistent pattern into a quantified one, and it should be done before the result is restated anywhere as a significant difference.

\textbf{Small-scale human calibration.} The clearest methodological gap is the missing human ground truth (Section~\ref{sec:limitations}). The tractable next step is not annotating everything but building a small expert-annotated anchor set, say 200 articles scored by an oil-market professional, to benchmark both the LLM extraction and the cross-model agreement against. That would turn the inter-model calibration from a standalone signal into a calibrated one, without paying to annotate the full corpus.

\textbf{Separating tone and price sentiment as parallel features.} The construct-dependence finding (Section~\ref{sec:rq2-discussion}) suggests an experiment, not just a caveat: carry the FinBERT tone score and the LLM price-impact score as separate features and let the model weigh them itself. For a liquidity target especially, where salience and risk drive volume more than price direction does, the tone signal may hold information the price signal does not. Their disagreement, the geopolitical sign-flip, might even work as a feature in its own right, marking high-uncertainty news.

\textbf{Multi-regime training.} The \texttt{price\_range} extrapolation failure (Section~\ref{sec:limitations}) is a training-window problem rather than a model-capacity one. As more regime episodes accumulate in the data, retraining across multiple structural breaks, or conditioning on a regime indicator explicitly, would test whether the volatility target becomes learnable once the model has seen more than one regime.

\textbf{Cross-commodity replication.} The risk-and-salience reading (Section~\ref{sec:rq2-discussion}) is a hypothesis about oil specifically, where geopolitical supply risk is price-bullish. Running the two-phase design on commodities with different supply and demand structures (natural gas, gold, copper) would test whether the same risk-premium mechanism drives their news-liquidity response or whether it belongs to WTI's geopolitical exposure. Cross-commodity spillovers, where oil news moves gas or gold liquidity, are a bigger question this thesis leaves open.

\textbf{Richer entity representation.} The entity flags are a binary vocabulary of 71 canonical actors (Section~\ref{sec:llm-extraction}). Swapping them for domain-aware entity embeddings, learned representations that would put Oman and Qatar near each other as Hormuz-adjacent producers, would let the model generalize across related actors instead of treating each one as an independent indicator. That could sharpen the risk-and-salience signal the current flags only approximate.
